%% kapa: exam=1 round=33 year=2022 period=1 form=A subject=civil_law subject_ko=민법 seq=37 number=37
%% source: Q-Net 감정평가사 기출문제 (cst00309) artlSeq=5209997 file=raw/1008/5209997/2022년도 제33회 감정평가사 1차시험 1교시(민법,경제학,부동산학).pdf
%% conversion: kordoc 4.24.51 -> tools/gen1.py (auto); review=proofread-sonnet; flags=none
\begin{question}{37}
민법 제365조의 일괄경매청구권에 관한 설명으로 옳은 것을 모두 고른 것은? (다툼이 있으면 판례에 따름)
\begin{framedtext}
ㄱ. 토지에 저당권을 설정한 후 그 설정자가 그 토지에 건물을 축조하여 저당권자가 토지와 함께 그 건물에 대하여도 경매를 청구하는 경우, 저당권자는 그 건물의 경매대가에 대해서도 우선변제를 받을 권리가 있다. \par
ㄴ. 저당권설정자로부터 저당토지에 대한 용익권을 설정받은 자가 그 토지에 건물을 축조한 후 저당권설정자가 그 건물의 소유권을 취득한 경우, 저당권자는 토지와 건물을 일괄하여 경매를 청구할 수 있다. \par
ㄷ. 토지에 저당권을 설정한 후 그 설정자가 그 토지에 축조한 건물의 소유권이 제3자에게 이전된 경우, 저당권자는 토지와 건물을 일괄하여 경매를 청구할 수 없다.
\end{framedtext}
\begin{choices}
  \item ㄱ
  \item ㄴ
  \item ㄷ
  \item ㄴ, ㄷ
  \item ㄱ, ㄴ, ㄷ
\end{choices}
\end{question}
\answer{4}
